\exposetitle{XXV}{The existence theorem}
\label{XXV}

\setcounter{footnote}{-1}
\nde{Version of 13/10/2024.}
For completeness, we give in this expos\'e a proof of the existence theorem for split groups. Like Chevalley's original proof (``Sur certains sch\'emas de groupes semi-simples'', S\'eminaire Bourbaki, May 1961, no.~219), it rests on the existence of complex semisimple algebraic groups of all possible types. The principle of a more satisfactory proof, proving directly the existence of a simply connected split semisimple $\ZZ$-group corresponding to a given Cartan matrix, was given by Cartier (unpublished).\nde{Such a principle was also outlined by B.~Kostant \cite{Ko66}; a complete proof, using quantum groups, was given recently by G.~Lusztig \cite{Lu09}.} We point out, however, that the difficulty is not to give an explicit construction of a group scheme, but to verify that the group thus constructed does satisfy the required conditions, that is, essentially, that its fibers are indeed smooth and reductive.

\sgasection{1}{Statement of the theorem}
\begin{theorem}{1.1}\label{XXV.1.1}
Let $S$ be a nonempty scheme. The functor
\[
G\mto\mathscr R(G)
\]
is an equivalence from the category of pinned reductive $S$-groups to the category of pinned reduced root data.
\end{theorem}
By the uniqueness theorem (Exp.~XXIII, 4.1), the preceding statement is equivalent to:
\begin{corollary}{1.2 (Existence of split groups)}\label{XXV.1.2}
For every reduced root datum $\mathscr R$, there exists a split reductive $\ZZ$-group $G$ such that $\mathscr R(G)\simeq\mathscr R$.\nde{We also point out that the work \cite{BT84} of F.~Bruhat and J.~Tits contains a variant of Chevalley's construction (\loccit, 2.2.3--2.2.5 and \S~3.2), which provides in particular a smooth affine $\ZZ$-group $G$ with connected fibers, possessing a split maximal torus, and whose generic fiber is a reductive $\QQ$-group of type $\mathscr R$; the fact that the geometric fibers of $G$ are reductive follows from the study of the unipotent radical of a special fiber made in \loccit, 4.6.12 and 4.6.15 (valid for more general $G$, associated with a valued root datum), but it is simpler to deduce it from the description of $\mathfrak g=\Lie(G)$, from Exp.~XIX 1.12~(iii), and from the existence of the elements $w_\alpha(X)$ of Exp.~XX 3.1~(iv) (cf.~\cite{BT84}, 3.2.1).}
\end{corollary}

% Source PDF p. 2.
In particular:
\begin{corollary}{1.3}\label{XXV.1.3}
Let $k$ be a field. For every splittable reductive $k$-group $G_k$, there exists a reductive $\ZZ$-group $G$ such that $G\otimes_\ZZ k\simeq G_k$.
\end{corollary}
Conversely, let us first note that, to prove 1.2, it suffices by Exp.~XXII 4.3.1 and Exp.~XXI 6.5.10 to consider the case where the root datum $\mathscr R$ is simply connected (and even irreducible, if one wishes, by Exp.~XXI, 7.1.6).

Moreover, under the conditions of 1.3, the group scheme $G$ has constant type (since $\Spec(\ZZ)$ is connected), hence type $\mathscr R(G_k)$; by Exp.~XXIII, 5.9, it follows that the validity of 1.3 for a given group $G_k$ implies the existence of a split $\ZZ$-group of type $\mathscr R(G_k)$.\nde{In fact, Exp.~XXIII, 5.9 is not necessary, since the present expos\'e constructs, for every semisimple $\QQ$-group $G_\QQ$, a semisimple $\ZZ$-group $G$ of the same type as $G_\QQ$, equipped with a split maximal torus $T$; thus, by Exp.~XXII 2.2, $G$ is split.}

To prove 1.2 and hence 1.1, it therefore suffices to prove 1.3 when $k$ has characteristic zero (for example $k=\CC$) and $G_k$ is simply connected (and in particular semisimple), as well as:
\begin{proposition}{1.4}\label{XXV.1.4}
For every simply connected reduced root datum $\mathscr R$, there exists a semisimple algebraic $\CC$-group of type $\mathscr R$.
\end{proposition}
One can prove 1.4 as follows. First, one knows that there exists a complex semisimple Lie algebra $\mathfrak g$ of type $\mathscr R$, cf.~for example N.~Jacobson, \textit{Lie Algebras}, ch.~VII, Th.~5.\nde{Let $\Delta$ be a base of $R$, and let $\widetilde{\mathfrak g}$ be the $\CC$-Lie algebra generated by generators $(e_\alpha,f_\alpha,h_\alpha)_{\alpha\in\Delta}$ subject to the relations $[h_\alpha,h_\beta]=0$, $[e_\alpha,f_\beta]=h_\alpha$ if $\beta=\alpha$ and $=0$ otherwise, $[h_\alpha,e_\beta]=(\alpha^*,\beta)e_\beta$ and $[h_\alpha,f_\beta]=-(\alpha^*,\beta)f_\beta$. In \loccit, $\mathfrak g$ is defined as the quotient of $\widetilde{\mathfrak g}$ by the intersection of the kernels of the finite-dimensional irreducible representations of $\widetilde{\mathfrak g}$; in \cite{Se66}, \S~VI.5, Th.~9 (see also \cite{BLie}, VIII, \S~4.3, Th.~1), it is shown that $\mathfrak g$ is the quotient of $\widetilde{\mathfrak g}$ by the relations $\ad(e_\alpha)^{1-(\alpha^*,\beta)}(e_\beta)=0$ and $\ad(f_\alpha)^{1-(\alpha^*,\beta)}(f_\beta)=0$. For an explicit description of the structure constants (in particular the choice of signs), see Exp.~XXIII 6.5 and 6.7, as well as \cite{Ti66}, \S~4, Th.~1, and (for types A,D,E) \cite{Sp98}, 10.2.5.} Then $G=\Aut(\mathfrak g)^0$ is a semisimple algebraic $\CC$-group of type $\ad(\mathscr R)$.\nde{One can suppose $\mathscr R$ irreducible, hence $\mathfrak g$ simple. By a general argument, one knows that $\Lie(G)$ is the $\CC$-Lie algebra of derivations of $\mathfrak g$ (cf.~\cite{DG70}, \S~II.4, Prop.~2.3); these are all inner (cf.~\cite{BLie}, I \S~6.1, cor.~3 of prop.~1), so $\Lie(G)=\mathfrak g$; consequently $G$ has no invariant subgroup of dimension $>0$, hence $G$ is semisimple. Its root system is then the same as that of $\mathfrak g$; moreover the center of $G$ acts both trivially and faithfully on $\mathfrak g$, and so is trivial, hence $G$ is adjoint.} By Bible, \S~23.1, prop.~1, one deduces the existence of a semisimple $\CC$-group of type $\mathscr R$.

The remainder of this expos\'e is devoted to proving 1.3 for $k$ of characteristic zero and $G_k$ semisimple. This will be done in two stages: construction of a ``chunk of a group $\ZZ$-scheme'' (no.~2), study of the group obtained by applying the ``Weil theorem'' (no.~3).

% Source PDF p. 3.
To avoid confusion, \emph{we shall not use in no.~2} the abbreviated notation $X_k$ to denote the $k$-scheme $X\otimes_\ZZ k$, where $X$ is a $\ZZ$-scheme.

\sgasection{2}{Existence theorem: construction of a group chunk}
\numpara{2.1}\label{XXV.2.1}
Let us choose once and for all a splitting of $G_k$, denoted
\[
(G_k,T_k,M,R)
\]
(cf.~Exp.~XXII, 1.13), a system of simple roots $\Delta$ of $R$ (defining the system of positive roots $R_+$), and a Chevalley system $(X_{\alpha,k})_{\alpha\in R}$ of $G_k$ (Exp.~XXIII, 6.1 and 6.2) satisfying the following additional condition (cf.~XX 2.6): for every $\alpha\in R$, one has
\[
X_{\alpha,k}X_{-\alpha,k}=1.
\]
Finally, let us choose on the subgroup of $M$ generated by $R$ a total ordering compatible with the group structure, such that the roots $>0$ are the elements of $R_+$. The roots are then denoted
\[
-\alpha_n<-\alpha_{n-1}<\cdots<-\alpha_1<\alpha_1<\alpha_2<\cdots<\alpha_n.
\]
For $\alpha\in R$, one denotes by $U_{\alpha,k}$ the vector group corresponding to the root $\alpha$, and by
\[
p_{\alpha,k}:\mathbf G_{\mathrm a,k}\isomto U_{\alpha,k}
\]
the isomorphism of vector groups defined by $X_{\alpha,k}$.

\numpara{2.2}\label{XXV.2.2}
The splitting of $G_k$ includes in particular an isomorphism of $k$-groups
\[
T_k\simeq D_k(M).
\]
Put $T=D(M)$; this is a $\ZZ$-torus, and one can regard the preceding isomorphism as an isomorphism
\[
T_k\simeq T\otimes_\ZZ k.
\]
One has
\[
\Hom_{\ZZ\text{-gr.}}(T,\Gm)=M,
\]
and one will regard the elements of $R\subset M$ as characters of $T$. Similarly, one will regard the elements of $R^*$ as morphisms of $\ZZ$-groups $\Gm\to T$.

\numpara{2.3}\label{XXV.2.3}
For each $\alpha\in R_+$, let $\Ga(\alpha)$ be a copy of the group $\Ga$; consider the $\ZZ$-scheme
\[
U=\Ga(\alpha_1)\times\cdots\times\Ga(\alpha_n).
\]
If $U_k$ denotes the unipotent part of the Borel group $B_k$ of $G_k$ defined by $R_+$, denote by
\[
a:U\otimes_\ZZ k\isomto U_k
\]
the isomorphism of $k$-schemes defined by
\[
a(x_1,\ldots,x_n)=p_{\alpha_1,k}(x_1)\cdots p_{\alpha_n,k}(x_n).
\]

% Source PDF p. 4.
\numpara{2.4}\label{XXV.2.4}
The group law of $U_k$ is expressed by relations of the form
\[
a(x_1,\ldots,x_n)\cdot a(y_1,\ldots,y_n)=a(z_1,\ldots,z_n),
\]
where each $z_h$ ($h=1,\ldots,n$) is expressed as a polynomial
\[
z_h=x_h+y_h+Q_h(x_1,\ldots,x_{h-1},y_1,\ldots,y_{h-1}),
\]
the coefficients of $Q_h$ being integers (Exp.~XXII, 5.5.8 and Exp.~XXIII, 6.4). Moreover $Q_h(x_1,\ldots,x_{h-1},0,\ldots,0)=0$.

Let us equip $U$ with the composition law defined by the preceding formulas (which are indeed ``defined over $\ZZ$''). Since this law induces on $U_k$ its group law, it is associative, and $(0)$ is an identity element for it (indeed, the preceding two assertions are expressed by relations among the polynomials $Q_h$, and $\ZZ\to k$ is injective). Let us show that it is a group law: if $(x_i)$ is a section of $U$ (over any $S$), one computes the inverse $(y_i)$ of $(x_i)$ by the recurrence formulas:
% typo?: first argument x_i in Q_i is kept as printed.
\[
y_i=-x_i-Q_i(x_i,\ldots,x_{i-1},y_1,\ldots,y_{i-1})
\]
which are again ``defined over $\ZZ$''.

To sum up, we have constructed on $U$ a group law such that the preceding isomorphism $a$ is an isomorphism of groups.

For each $\alpha\in R_+$, consider the morphism
\[
p_\alpha:\Ga\to U
\]
defined by $p_\alpha(x)=(x_i)$, where
\[
\begin{cases}
x_i=x&\text{if }\alpha_i=\alpha,\\
x_i=0&\text{if }\alpha_i\ne\alpha.
\end{cases}
\]
This is a closed immersion and a group homomorphism; its image is denoted $U_\alpha$. One has $(x_i)=p_{\alpha_1}(x_1)\cdots p_{\alpha_n}(x_n)$, which proves that $U$ is identified with the product
\[
U=U_{\alpha_1}\cdot U_{\alpha_2}\cdots U_{\alpha_n}.
\]

\numpara{2.5}\label{XXV.2.5}
Let us make $T=D(M)$ act on each $U_\alpha$ through the character $\alpha$; one immediately verifies that this defines an action of $T$ on the group $U$, and one can construct the semidirect product $B=T\cdot U$. One has a canonical isomorphism of $k$-groups
\[
B\otimes_\ZZ k\isomto B_k.
\]
If we now take any ordering on $R_+$, the morphism
\[
\prod_{\alpha\in R_+}U_\alpha\to U
\]
defined by the product in $U$ is still an isomorphism. Indeed, since both sides are flat $\ZZ$-schemes of finite presentation, one can content oneself with verifying the assertion on the geometric fibers; one is then reduced to Lazard's theory (Bible, \S~13.1, prop.~1): one regards $U$ as a group with operators $T$, and uses the fact that the $U_\alpha$ are pairwise nonisomorphic as groups with operators (since the characters $\alpha\in R_+$ of $T$ are pairwise distinct on each fiber).

% Source PDF p. 5.
\numpara{2.6}\label{XXV.2.6}
Replacing $R_+$ by $R_-=-R_+$, one similarly constructs groups $U^-$, $B^-$, $U_\alpha$ ($\alpha\in R_-$), and isomorphisms
\[
p_\alpha:\Ga\isomto U_\alpha,\qquad\alpha\in R_-.
\]
Finally, introduce the product scheme
\[
\Omega=U^-\times T\times U;
\]
one has a canonical isomorphism of $k$-schemes
\[
\Omega\otimes_\ZZ k\isomto U_k^-\times_k T_k\times_k U_k\simeq\Omega_k,
\]
where $\Omega_k$ is the ``big cell'' of $G_k$ (Exp.~XXII, 4.1.2).

From now on, we identify $\Omega\otimes_\ZZ k$ with $\Omega_k$ by the preceding isomorphism; we regard $U^-$, $T$, $U$ as subschemes of $\Omega$, by means of the identity sections. One denotes by $\underline e=((0),e,(0))$ the ``identity section'' of $\Omega$.

Our aim now is to put a group-chunk law on $\Omega$.
\begin{lemma}{2.7}\label{XXV.2.7}
Let $\alpha\in\Delta$, and let $w_{\alpha,k}$ be the element of $\uNorm_{G_k}(T_k)(k)$ defined by $X_{\alpha,k}$ (recall that by definition one has
\[
w_{\alpha,k}=p_{-\alpha,k}(-1)p_{\alpha,k}(1)p_{-\alpha,k}(-1)).
\]
There exists an open subset $V_\alpha$ of $\Omega$, containing the section $\underline e$, and a morphism
\[
h_\alpha:V_\alpha\to\Omega,
\]
satisfying the following conditions:
\begin{enumeratei}
\item $h_\alpha(\underline e)=\underline e$,
\item $(h_\alpha)\otimes_\ZZ k$ coincides with the restriction of $\operatorname{int}(w_{\alpha,k})$ to $V_\alpha\otimes_\ZZ k\subset G_k$.
\item One has $T\subset V_\alpha$, and $h_\alpha$ sends $T$ into $T$. For every $\beta\in R$, one has $U_\beta\subset V_\alpha$, and $h_\alpha$ sends $U_\beta$ into $U_{s_\alpha(\beta)}$.
\end{enumeratei}
\end{lemma}
By the definition of a Chevalley system (Exp.~XXIII, 6.1), there exists for each $\beta\in R$ an integer $e_\beta=\pm1$ such that
\[
\operatorname{int}(w_{\alpha,k})p_{\beta,k}(x)=p_{s_\alpha(\beta),k}(e_\beta x)
\]
for every $x\in\Ga(S)$, $S\to\Spec(k)$.

Let $S$ be any scheme; let us write any element of $\Omega(S)$ in the form
\[
\begin{split}
u=\biggl(&\Bigl(\prod_{\substack{\beta\in R_-\\\beta\ne-\alpha}}p_\beta(x_\beta)\Bigr)\cdot p_{-\alpha}(x_{-\alpha}),\quad t,\\
& p_\alpha(x_\alpha)\cdot\Bigl(\prod_{\substack{\beta\in R_+\\\beta\ne\alpha}}p_\beta(x_\beta)\Bigr)\biggr),
\end{split}
\]
where one has chosen an ordering (any ordering) on $R_--\{-\alpha\}$ and $R_+-\{\alpha\}$ (cf.~2.5).

One defines a morphism $d:\Omega\to\Spec(\ZZ)$ by
% typo?: the target Spec(Z) is kept as printed.
\[
d(u)=\alpha(t)+x_\alpha x_{-\alpha}.
\]

% Source PDF p. 6.
Let $V_\alpha$ be the open subset $\Omega_d$ (that is, the open subset of $\Omega$ defined by ``$d(u)$ invertible''); it contains $\underline e$, $T$, and each $U_\beta$, $\beta\in R$. Let
\[
h_\alpha:V_\alpha\to\Omega
\]
be the morphism defined by $h_\alpha(u)=(a(u),b(u),c(u))$, where
\[
\begin{aligned}
a(u)&=\Bigl(\prod_{\substack{\beta\in R_-\\\beta\ne-\alpha}}p_{s_\alpha(\beta)}(e_\beta x_\beta)\Bigr)\cdot p_{-\alpha}\bigl(-x_\alpha d(u)^{-1}\bigr),\\
b(u)&=t\cdot\alpha^*(d(u)),\\
c(u)&=p_\alpha\bigl(-x_{-\alpha}d(u)^{-1}\bigr)\cdot\Bigl(\prod_{\substack{\beta\in R_+\\\beta\ne\alpha}}p_{s_\alpha(\beta)}(e_\beta x_\beta)\Bigr).
\end{aligned}
\]
Since $s_\alpha$ permutes the positive (resp.~negative) roots distinct from $\alpha$ (resp.~$-\alpha$), $c(u)$ (resp.~$a(u)$) is then a section of $U$ (resp.~$U^-$), and the preceding morphism is well defined. It trivially satisfies (i) and (iii). As for (ii), this follows immediately from the definition of the $e_\beta$, $\beta\in R$, and from Exp.~XX, 3.12.

\begin{lemma}{2.8}\label{XXV.2.8}
There exist open subsets $V$ and $V'$ of $\Omega$, and morphisms
\[
h:V\to\Omega,\qquad h':V'\to\Omega,
\]
satisfying the following conditions:
\begin{enumeratei}
\item $V$ and $V'$ contain $\underline e$, and $h(\underline e)=h'(\underline e)=\underline e$.
\item The morphism induced by $h'\circ h:h^{-1}(V')\to\Omega$ is the restriction of the identity morphism.
\item The $k$-morphisms $h\otimes_\ZZ k$ and $h'\otimes_\ZZ k$ are the restrictions to $V\otimes_\ZZ k$ and $V'\otimes_\ZZ k$ of automorphisms of the group $G_k$.
\item $V$ and $V'$ contain $U$, $T$ and $U^-$; $h$ and $h'$ send $U$ into $U^-$, $U^-$ into $U$, and $T$ into $T$.
\end{enumeratei}
\end{lemma}
Let $\overline w_0$ be the element of the Weyl group of $G_k$ which transforms $R_+$ into $R_-$. Write
\[
\overline w_0=s_{\alpha_n}\cdots s_{\alpha_1},\qquad\alpha_i\in\Delta
\]
(no connection with the numbering of 2.1). Put
\[
w_0=w_{\alpha_n,k}\cdots w_{\alpha_1,k}\in\uNorm_{G_k}(T_k)(k).
\]
Define by induction on $i\le n$ an open subset $V_i$ of $\Omega$ and a morphism $h_i:V_i\to\Omega$ by $V_0=\Omega$, $h_0=\mathrm{id}$, and, for $i=0,\ldots,n-1$,
\[
V_{i+1}=h_i^{-1}(V_{\alpha_{i+1}}),\qquad h_{i+1}=h_{\alpha_{i+1}}\circ h_i,
\]
where the notations $V_{\alpha_j}$ and $h_{\alpha_j}$ are those of 2.7.

% Source PDF p. 7.
Take $V=V_n$ and $h=h_n$. The conditions of (i), (iii) and (iv) concerning $V$ and $h$ are indeed satisfied; for (i) and (iii) this follows immediately from 2.8, and for (iv), from the fact that $h\otimes_\ZZ k$ is the restriction of $\operatorname{int}(w_0)$ to $V\otimes_\ZZ k$.
% typo?: reference 2.8 is kept as printed.

Since $(\overline w_0)^2=1$, one also has
\[
\overline w_0=s_{\alpha_1}\cdots s_{\alpha_n}=(s_{\alpha_1}^3)\cdots(s_{\alpha_n}^3).
\]
Putting
\[
w'_0=(w_{\alpha_1,k})^3\cdots(w_{\alpha_n,k})^3,
\]
and carrying out the same construction as above, one deduces an open subset $V'$ and a morphism $h'$ likewise satisfying (i), (iii), (iv). Moreover, $h'\otimes_\ZZ k$ is the restriction of $\operatorname{int}(w'_0)$ to $V'\otimes_\ZZ k$. But for each simple root $\alpha\in\Delta$, one has $(w_{\alpha,k})^4=e$ (cf.~Exp.~XX, 3.1), so $w'_0\cdot w_0=e$, which shows that $h'\circ h$ induces the identity morphism in a nonempty open subset of $\Omega\otimes_\ZZ k$. But since $\Omega$ is smooth and of finite presentation over $\ZZ$, $\Omega\otimes_\ZZ k$ is schematically dense in $\Omega$, which proves (ii).

\begin{proposition}{2.9}\label{XXV.2.9}
There exist an open subset $V_1$ of $\Omega\times\Omega$, an open subset $V_2$ of $\Omega$, and morphisms
\[
\pi:V_1\to\Omega,\qquad\sigma:V_2\to\Omega,
\]
possessing the following properties:
\begin{enumeratei}
\item If $x\in\Omega(S)$, then $(\underline e,x)$ and $(x,\underline e)$ are sections of $V_1$, and
\[
\pi(\underline e,x)=\pi(x,\underline e)=x.
\]
\item $V_2$ contains $\underline e$, and $\sigma(\underline e)=\underline e$.
\item $\pi_k$ and $\sigma_k$ are the restrictions of the morphisms $G_k\times_k G_k\to G_k$ and $G_k\to G_k$ defined by the product (resp.~the inverse).
\end{enumeratei}
\end{proposition}
\begin{proof}
Let $(v,t,u)$ and $(v',t',u')$ be two sections of $\Omega$. Then $h(u)$ is a section of $U^-$, $h(v')$ is a section of $U$ by 2.8~(iv), and one can therefore regard $(h(u),e,h(v'))$ as a section of $\Omega$. Let $V_1$ be the open subset of $\Omega\times\Omega$ defined by the condition:
\[
\bigl((v,t,u),(v',t',u')\bigr)\in V_1(S)\ \Longleftrightarrow\ (h(u),e,h(v'))\in V'(S)
\]
(notations of 2.8). If $\bigl((v,t,u),(v',t',u')\bigr)$ is a section of $V_1$, then $h'(h(u),e,h(v'))$ is defined; this is a section of $\Omega$ which one can decompose:
% typo: h'(h(u),e,h(v'))) -> h'(h(u),e,h(v')).
\[
h'(h(u),e,h(v'))=(v'',t'',u'').
\]
One then puts
\[
\pi\bigl((v,t,u),(v',t',u')\bigr)=(v\cdot tv''t^{-1},\ tt''t',\ t'^{-1}u''t'\cdot u').
\]
The verification of (i) is immediate (by 2.8~(ii)). To verify the condition of (iii) concerning $\pi$, one sees that
\[
h'(h(u),e,h(v'))=uv'=v''t''u''
\]
when $u\in U(S)$, $v\in U^-(S)$, $S\to k$, by 2.8~(iii) and (ii).
% typo?: v (rather than v') is kept as printed in the preceding condition.
One constructs $\sigma$ in a similar manner: if $(v,t,u)$ is a section of $\Omega$, $h(u^{-1})$ is a section of $U^-$, $h(v^{-1})$ a section of $U$, $h(t^{-1})$ a section of $T$, so $(h(u^{-1}),h(t^{-1}),h(v^{-1}))$ is a section of $\Omega$, and one can define an open subset $V_2$ of $\Omega$ by
\[
(v,t,u)\in V_2(S)\ \Longleftrightarrow\ (h(u^{-1}),h(t^{-1}),h(v^{-1}))\in V'(S)
\]
and a morphism $\sigma:V_2\to\Omega$ by
\[
\sigma(v,t,u)=h'(h(u^{-1}),h(t^{-1}),h(v^{-1})).
\]
% Source PDF p. 8: the preceding sentence began on p. 7.
One verifies the conditions on $\sigma$ as above.
\end{proof}
\begin{corollary}{2.10}\label{XXV.2.10}
$\pi$ is ``generically associative'', and $\sigma$ is a ``generic inverse'': if $x,y,z\in\Omega(S)$ and if the expressions below are defined (which always occurs above an open subset of $\Omega$ containing the identity section), one has:
\[
\pi(x,\pi(y,z))=\pi(\pi(x,y),z),\qquad
\pi(x,\sigma(x))=\underline e=\pi(\sigma(x),x).
\]
\end{corollary}
Indeed, the two sides of each of these formulas define morphisms between smooth $\ZZ$-schemes of finite presentation, which coincide on the generic fibers, by 2.10~(iii).
% typo?: reference 2.10 (iii) is kept as printed.
\begin{corollary}{2.11}\label{XXV.2.11}
Let $\alpha\in R$. For every $S$ and all $x,y\in\Ga(S)$ such that $(p_\alpha(x),p_{-\alpha}(y))\in V_1(S)$ and $1+xy\in\Gm(S)$ (which defines an open subset of $\mathbf G_{\mathrm a,S}^2$ containing the section $(0,0)$), one has:
\[
\pi(p_\alpha(x),p_{-\alpha}(y))=
\Bigl(p_{-\alpha}\Bigl(\frac{y}{1+xy}\Bigr),\ \alpha^*(1+xy),\ p_\alpha\Bigl(\frac{x}{1+xy}\Bigr)\Bigr).
\]
\end{corollary}
The proof is the same as before, by Exp.~XX, 2.1.

\sgasection{3}{Existence theorem: end of the proof}
To simplify the language, let us make the following definition.
\begin{definition}{3.1}\label{XXV.3.1}
Let $S$ be a scheme and $G$ a group $S$-scheme. One says that $G$ is \emph{admissible} if there exists an open immersion of $S$-schemes $i:\Omega_S=\Omega\times S\to G$ satisfying the following conditions:
\begin{enumeratei}
\item The diagram
\[
\xymatrix@C=18pt@R=18pt{
(V_1)_S\ar[r]\ar[d]_{\pi}&\Omega_S\times_S\Omega_S\ar[r]^{i\times_S i}&G\times_S G\ar[d]^{\pi_G}\\
\Omega_S\ar[rr]_{i_S}&&G,
}
\]
where $\pi_G$ denotes the multiplication morphism in $G$, is commutative.
\item There exists a finite set of sections $a_j\in\Omega(S)$ such that the $i(a_j)\cdot i(\Omega_S)$ cover $G$.
\end{enumeratei}
\end{definition}
By the ``Weil theorem'' (Exp.~XVIII, 3.13~(iii) and (iv)), one has:
\begin{lemma}{3.2}\label{XXV.3.2}
If, for every scheme $S$ \emph{\'etale and of finite type} over $\ZZ$, every admissible $S$-group is affine, then there exists an admissible affine $\ZZ$-group.
\end{lemma}
